% GENERATED FILE. Edit canonical lessons, then run npm run generate:books.
% canonical-source-hash: fnv1a64:6bd39a26f9cc590c
% canonical-lessons: ZH-C404-shoupiaochu, ZH-C404-jianpiao, ZH-W404-jian-check, ZH-C404-bupiao, ZH-C404-wopu, ZH-C404-yingzuo

\chapter{Tickets, Please}
\label{ch:zh-tickets-please}
\hlchaptermodality{404}

\begin{chapteropening}
\textbf{By the end of this chapter:} \emph{I can find the ticket office, have my ticket checked, pay on board and choose a berth or a seat.}

Complete the last lesson of chapter 404: I can find the ticket office, have my ticket checked, pay on board and choose a berth or a seat.
\end{chapteropening}

\section[shòupiàochù]{shòupiàochù — ticket office}

\label{lesson:ZH-C404-shoupiaochu}



\begin{quote}
Before the new one: say the Mandarin for a train attendant, then the Mandarin for the head of a station.
\end{quote}

\subsection*{You'll want to know: shòupiàochù}
\emph{shòupiàochù} — \textquotedblleft{}ticket office\textquotedblright{}.

\begin{quote}
\emph{qù shòupiàochù zěnme zǒu} — How do I get to the ticket office?
\end{quote}

Its characters come later; for now it is a word for the ear.

\begin{cousinweb}
\emph{shòupiàochù} is built from \emph{shòu} 'to sell', \zh{票} 'a ticket' and \emph{chù} 'a place'. The first and the last are inherited from Middle Chinese, \emph{dzyuwH} and \emph{tsyhoH} (Guangyun, in Baxter's transcription), glossed 'to sell goods' and 'a place'; no dictionary we use records an older form of \zh{票}, so its history is not settled.
\end{cousinweb}

\begin{grammarlens}[title={What you've built}]
One more word for getting around.
\end{grammarlens}

\subsection*{Your turn}
\begin{itemize}
\raggedright
  \item \emph{Say it:} \emph{shòupiàochù}
  \item \emph{Say it:} \emph{shòupiàochù}, once more
  \item \emph{Say it:} \emph{qù shòupiàochù zěnme zǒu} — How do I get to the ticket office?
  \item \emph{Recall:} say the Mandarin for a train attendant, then the Mandarin for the head of a station, then say \emph{shòupiàochù} again
\end{itemize}

\begin{tcolorbox}[breakable,colback=teal!4,colframe=teal!35!black,title={Before you move on}]
What does shòupiàochù mean? (\textquotedblleft{}Ticket office\textquotedblright{}.) Say it once more.
\end{tcolorbox}
\section[\texorpdfstring{\zh{检票} (jiǎnpiào)}{检票 (jiǎnpiào)}]{\zh{检票} (jiǎnpiào) — to check tickets}

\label{lesson:ZH-C404-jianpiao}



\begin{quote}
Before the new one: say the Mandarin for the head of a station, then the Mandarin for a ticket office.
\end{quote}

\subsection*{You'll want to know: \zh{检票}}
\textbf{\zh{检票}} — \emph{jiǎnpiào} — \textquotedblleft{}to check tickets\textquotedblright{}.

\begin{quote}
\textbf{\zh{现在检票了}} \emph{xiànzài jiǎnpiào le} — They're checking tickets now.
\end{quote}

The new character, in pieces: \zh{检} — eleven strokes in two pieces: on the left the tree piece, four strokes, as on the left of \zh{机}; on the right a long stroke falling to the left and one falling to the right, like a roof, then a short bar, two dots and a short stroke falling to the left, and a long bar across the foot. The next lesson puts it in your hand.

\begin{cousinweb}
\textbf{\zh{检票}} (\emph{jiǎnpiào}) is built from \zh{检} 'to check' and \zh{票} 'a ticket'. The first is inherited from Middle Chinese \emph{kjemX} (Guangyun, in Baxter's transcription), glossed as the seal on a document and 'to check'; no dictionary we use records an older form of \zh{票}, so its history is not settled.
\end{cousinweb}

\begin{grammarlens}[title={What you've built}]
One more word for getting around.
\end{grammarlens}

\subsection*{Your turn}
\begin{itemize}
\raggedright
  \item \emph{Say it:} \emph{jiǎnpiào}
  \item \emph{Say it:} \emph{jiǎnpiào}, once more
  \item \emph{Say it:} \emph{xiànzài jiǎnpiào le} — They're checking tickets now.
  \item \emph{Recall:} say the Mandarin for the head of a station, then the Mandarin for a ticket office, then say \emph{jiǎnpiào} again
\end{itemize}

\begin{tcolorbox}[breakable,colback=teal!4,colframe=teal!35!black,title={Before you move on}]
What does \zh{检票} mean? (\textquotedblleft{}To check tickets\textquotedblright{}.) Say it once more.
\end{tcolorbox}
\section[\texorpdfstring{\zh{检} (jiǎn)}{检 (jiǎn)}]{\zh{检} — the tree piece, then a roof and what is under it}

\label{lesson:ZH-W404-jian-check}



\begin{quote}
The word first, then the character inside it.

\begin{itemize}
\raggedright
  \item \emph{Recall:} say \emph{jiǎnpiào} — \textbf{R1}, one lesson back, before it decays
\end{itemize}

\begin{quote}
\zh{机} \zh{票}
\end{quote}

\emph{Read it:} each one aloud, with nothing said first — \emph{jī}, \emph{piào}
\end{quote}

\subsection*{Script — the tree piece, then a roof and what is under it}
\hlblockfigure{\detokenize{figures/ZH-W404-jian-check-filmstrip.pdf}}{How \zh{检} is written, stroke by stroke}

\begin{quote}
\zh{检}
\end{quote}

\textbf{Two pieces, left before right.} Count the pieces, not the strokes.

\textbf{The left piece, the tree piece, four strokes.} A short bar. A vertical down through it. A stroke falling to the left from the crossing. Then a short dot to the right.

\textbf{The right piece, seven strokes.} A long stroke falling to the left and a long stroke falling to the right: a roof. Under it a short bar. Then two dots side by side and a short stroke falling to the left beside them. Last, a long bar across the foot.

4 + 7 = eleven.

\subsection*{Your turn}
\begin{itemize}
\raggedright
  \item \emph{Write it:} \zh{检}, all eleven, with the model in front of you
  \item \emph{Cover:} write it again from memory, and count the pieces, not the strokes
\end{itemize}

\begin{tcolorbox}[breakable,colback=teal!4,colframe=teal!35!black,title={Before you move on}]
How many pieces does \zh{检} have? (\textbf{Two}.) Which comes first? (\textbf{The tree piece, as in \zh{机}}.) What closes the foot? (\textbf{A long bar}.)

Source: \href{https://raw.githubusercontent.com/chanind/hanzi-writer-data/68d10a4b21150cae5e1ebbd223eed289cf32d90c/data/\%E6\%A3\%80.json}{Hanzi Writer Data \zh{检}, pinned stroke paths and medians}.
\end{tcolorbox}
\section[bǔpiào]{bǔpiào — to buy a ticket on board}

\label{lesson:ZH-C404-bupiao}



\begin{quote}
Before the new one: say the Mandarin for a ticket office, then the Mandarin for to check tickets.
\end{quote}

\subsection*{You'll want to know: bǔpiào}
\emph{bǔpiào} — \textquotedblleft{}to buy a ticket on board\textquotedblright{}.

\begin{quote}
\emph{méiyǒu piào dehuà kěyǐ shàngchē bǔpiào} — If you've no ticket, you can buy one on the train.
\end{quote}

Its characters come later; for now it is a word for the ear.

\begin{cousinweb}
\emph{bǔpiào} is built from \emph{bǔ} 'to mend' and \zh{票} 'a ticket'. The first is inherited from Middle Chinese \emph{puX} (Guangyun, in Baxter's transcription), glossed 'to patch'; no dictionary we use records an older form of \zh{票}, so its history is not settled.
\end{cousinweb}

\begin{grammarlens}[title={What you've built}]
One more word for getting around.
\end{grammarlens}

\subsection*{Your turn}
\begin{itemize}
\raggedright
  \item \emph{Say it:} \emph{bǔpiào}
  \item \emph{Say it:} \emph{bǔpiào}, once more
  \item \emph{Say it:} \emph{méiyǒu piào dehuà kěyǐ shàngchē bǔpiào} — If you've no ticket, you can buy one on the train.
  \item \emph{Recall:} say the Mandarin for a ticket office, then the Mandarin for to check tickets, then say \emph{bǔpiào} again
\end{itemize}

\begin{tcolorbox}[breakable,colback=teal!4,colframe=teal!35!black,title={Before you move on}]
What does bǔpiào mean? (\textquotedblleft{}To buy a ticket on board\textquotedblright{}.) Say it once more.
\end{tcolorbox}
\section[wòpù]{wòpù — sleeping berth}

\label{lesson:ZH-C404-wopu}



\begin{quote}
Before the new one: say the Mandarin for to check tickets, then the Mandarin for to buy a ticket on board.
\end{quote}

\subsection*{You'll want to know: wòpù}
\emph{wòpù} — \textquotedblleft{}sleeping berth\textquotedblright{}.

\begin{quote}
\emph{wǒ xiǎng mǎi liǎng zhāng wòpù} — I'd like two sleeper berths.
\end{quote}

Its characters come later; for now it is a word for the ear.

\begin{cousinweb}
\emph{wòpù} is built from \emph{wò} 'to lie down' and \emph{pù} 'a bed', both inherited from Middle Chinese, \emph{ngwaH} and \emph{phuH} (Guangyun, in Baxter's transcription), which glosses the first 'to sleep' and the second 'to set out, to lay'; the bed sense of the second is not given there.
\end{cousinweb}

\begin{grammarlens}[title={What you've built}]
One more word for getting around.
\end{grammarlens}

\subsection*{Your turn}
\begin{itemize}
\raggedright
  \item \emph{Say it:} \emph{wòpù}
  \item \emph{Say it:} \emph{wòpù}, once more
  \item \emph{Say it:} \emph{wǒ xiǎng mǎi liǎng zhāng wòpù} — I'd like two sleeper berths.
  \item \emph{Recall:} say the Mandarin for to check tickets, then the Mandarin for to buy a ticket on board, then say \emph{wòpù} again
\end{itemize}

\begin{tcolorbox}[breakable,colback=teal!4,colframe=teal!35!black,title={Before you move on}]
What does wòpù mean? (\textquotedblleft{}Sleeping berth\textquotedblright{}.) Say it once more.
\end{tcolorbox}
\section[yìngzuò]{yìngzuò — hard seat}

\label{lesson:ZH-C404-yingzuo}



\begin{quote}
Before the new one: say the Mandarin for to buy a ticket on board, then the Mandarin for a sleeping berth.
\end{quote}

\subsection*{You'll want to know: yìngzuò}
\emph{yìngzuò} — \textquotedblleft{}hard seat\textquotedblright{}.

\begin{quote}
\emph{wǒ zuò yìngzuò qù Běijīng} — I'm going to Beijing in a hard seat.
\end{quote}

Its characters come later; for now it is a word for the ear.

\begin{cousinweb}
\emph{yìngzuò} is built from \emph{yìng} 'hard' and \emph{zuò} 'a seat', both inherited from Middle Chinese, \emph{ngeangH} and \emph{dzwaH} (Guangyun, in Baxter's transcription), which glosses the first 'firm' and the second as the seat of a bed.
\end{cousinweb}

\begin{grammarlens}[title={What you've built}]
One more word for getting around.
\end{grammarlens}

\subsection*{Your turn}
\begin{itemize}
\raggedright
  \item \emph{Say it:} \emph{yìngzuò}
  \item \emph{Say it:} \emph{yìngzuò}, once more
  \item \emph{Say it:} \emph{wǒ zuò yìngzuò qù Běijīng} — I'm going to Beijing in a hard seat.
  \item \emph{Recall:} say the Mandarin for to buy a ticket on board, then the Mandarin for a sleeping berth, then say \emph{yìngzuò} again
\end{itemize}

\begin{tcolorbox}[breakable,colback=teal!4,colframe=teal!35!black,title={Before you move on}]
What does yìngzuò mean? (\textquotedblleft{}Hard seat\textquotedblright{}.) Say it once more.
\end{tcolorbox}
